% Преамбула: пакеты, геометрия, шрифты, оформление заголовков.
% Требования см. в requirements.md и CLAUDE.md.
% Движок: LuaLaTeX.

% --- Язык ------------------------------------------------------------------
\usepackage{polyglossia}
\setdefaultlanguage{russian}
\setotherlanguage{english}

% --- Шрифты ----------------------------------------------------------------
% CMU Serif — Computer Modern Unicode (cm-unicode), кириллица + латиница.
% Соответствует требованию CLAUDE.md: «Times или Computer Modern».
% Пакет: texlive-fontsextra (Arch).
\usepackage{fontspec}
\defaultfontfeatures{Ligatures=TeX}
\setmainfont{CMU Serif}
\setsansfont{CMU Sans Serif}
\setmonofont{CMU Typewriter Text}
\newfontfamily\cyrillicfont{CMU Serif}
\newfontfamily\cyrillicfontsf{CMU Sans Serif}
\newfontfamily\cyrillicfonttt{CMU Typewriter Text}

% --- Геометрия страницы ----------------------------------------------------
\usepackage[a4paper, top=2cm, bottom=2cm, left=3cm, right=1cm]{geometry}

% --- Межстрочный интервал --------------------------------------------------
\usepackage{setspace}
\onehalfspacing

% --- Абзацный отступ -------------------------------------------------------
\usepackage{indentfirst}
\setlength{\parindent}{1.25cm}
\setlength{\parskip}{0pt}

% --- Графика, таблицы, математика ------------------------------------------
\usepackage{graphicx}
\graphicspath{{images/}}
\usepackage{amsmath, amssymb}
\usepackage{array}
\usepackage{booktabs}
\usepackage{longtable}
\usepackage{caption}

% --- Подписи рисунков и таблиц по требованиям ------------------------------
% Рисунок: подпись снизу по центру, формат «Рис. N. Название».
% Таблица: подпись сверху справа, шрифт 12pt, формат «Таблица N. Название».
\captionsetup[figure]{
  name=Рис.,
  labelsep=period,
  justification=centering,
  singlelinecheck=true,
  position=below
}
\captionsetup[table]{
  name=Таблица,
  labelsep=period,
  justification=raggedleft,
  singlelinecheck=false,
  font=small,
  position=above
}

% --- Нумерация страниц: внизу по центру, сквозная --------------------------
\usepackage{fancyhdr}
\pagestyle{fancy}
\fancyhf{}
\fancyfoot[C]{\thepage}
\renewcommand{\headrulewidth}{0pt}
\renewcommand{\footrulewidth}{0pt}

% Стиль для титульного листа: без номера, но страница считается первой.
\fancypagestyle{thesistitle}{%
  \fancyhf{}%
  \renewcommand{\headrulewidth}{0pt}%
  \renewcommand{\footrulewidth}{0pt}%
}

% --- Оформление заголовков (требование: строчные с первой прописной, ------
% полужирные, по левому краю с абзацным отступом, без точки в конце) -------
% Главы (\section): 16pt; разделы и подразделы: 14pt; всё полужирное,
% по левому краю с абзацным отступом, без точки в конце.
\usepackage{titlesec}
\titleformat{\section}
  {\normalfont\fontsize{16pt}{19pt}\bfseries}
  {\thesection.}{0.5em}{}
\titleformat{\subsection}
  {\normalfont\fontsize{14pt}{17pt}\bfseries}
  {\thesubsection.}{0.5em}{}
\titleformat{\subsubsection}
  {\normalfont\fontsize{14pt}{17pt}\bfseries}
  {\thesubsubsection.}{0.5em}{}

% Введение/Заключение/Список литературы/Приложения/Содержание — без номера,
% с новой страницы, с абзацным отступом, по левому краю, размер главы.
\newcommand{\thesisChapterStar}[1]{%
  \clearpage
  \phantomsection
  \addcontentsline{toc}{section}{#1}%
  \noindent\hspace*{\parindent}\textbf{\fontsize{16pt}{19pt}\selectfont #1}\par
  \vspace{\baselineskip}%
}

% --- Содержание ------------------------------------------------------------
% Стиль вдохновлён vcourse-document-template:
% секция — полужирные название и номер страницы, подсекция — полужирный
% номер, обычные название и страница; единые точечные «лидеры» на всех
% уровнях. Заголовок «Содержание» — по требованиям ВКР: с прописной буквы,
% по левому краю с абзацным отступом, полужирный, без точки в конце.
\usepackage{tocloft}

\renewcommand{\cftsecfont}{\bfseries}
\renewcommand{\cftsecpagefont}{\bfseries}

\renewcommand{\cftsubsecfont}{\mdseries}
\renewcommand{\cftsubsecpresnum}{\bfseries}
\renewcommand{\cftsubsecaftersnum}{\mdseries}
\renewcommand{\cftsubsecpagefont}{\mdseries}

\renewcommand{\cftdotsep}{2}
\renewcommand{\cftsecdotsep}{\cftdotsep}
\renewcommand{\cftsecleader}{\cftdotfill{\cftsecdotsep}}

\setlength{\cftbeforesecskip}{4pt}

\renewcommand{\contentsname}{Содержание}

% Печатаем заголовок «Содержание» сами, в стиле \thesisChapterStar,
% а затем зовём \@starttoc, чтобы tocloft не выводил собственный заголовок.
\makeatletter
\newcommand{\makethesistoc}{%
  \clearpage
  \thispagestyle{fancy}%
  \noindent\hspace*{\parindent}\textbf{\fontsize{16pt}{19pt}\selectfont \contentsname}\par
  \vspace{\baselineskip}%
  \@starttoc{toc}%
  \clearpage
}
\makeatother

% --- Гиперссылки (последним) ----------------------------------------------
\usepackage[hidelinks, unicode]{hyperref}

% --- API метаданных (значения подставляются в meta.tex) -------------------
\newcommand{\thesisMinistry}[1]{\gdef\thethesisMinistry{#1}}
\newcommand{\thesisInstitutionLong}[1]{\gdef\thethesisInstitutionLong{#1}}
\newcommand{\thesisUniversity}[1]{\gdef\thethesisUniversity{#1}}
\newcommand{\thesisDepartment}[1]{\gdef\thethesisDepartment{#1}}
\newcommand{\thesisTitle}[1]{\gdef\thethesisTitle{#1}}
\newcommand{\thesisProgramKind}[1]{\gdef\thethesisProgramKind{#1}}
\newcommand{\thesisProgramCode}[1]{\gdef\thethesisProgramCode{#1}}
\newcommand{\thesisProgramName}[1]{\gdef\thethesisProgramName{#1}}
\newcommand{\thesisHeadName}[1]{\gdef\thethesisHeadName{#1}}
\newcommand{\thesisAdvisorRank}[1]{\gdef\thethesisAdvisorRank{#1}}
\newcommand{\thesisAdvisorName}[1]{\gdef\thethesisAdvisorName{#1}}
\newcommand{\thesisStudentGroup}[1]{\gdef\thethesisStudentGroup{#1}}
\newcommand{\thesisStudentName}[1]{\gdef\thethesisStudentName{#1}}
\newcommand{\thesisCity}[1]{\gdef\thethesisCity{#1}}
\newcommand{\thesisYear}[1]{\gdef\thethesisYear{#1}}
